\section{Preliminaries}
\label{sec:prelim}

\paragraph{Analysis problem.}
A Python web handler is vulnerable to path traversal if attacker-controlled data can make a filesystem operation act on a path outside the directory the application intends to expose. Such operations include opening, reading, writing, deleting, listing, or probing existence or metadata. The attacker may use \texttt{..} components, absolute paths (which make \texttt{os.path.join} discard its prefix), encodings decoded after a check, or prefix confusion (\texttt{/srv/files2} passes a string check for \texttt{/srv/files}). We assume POSIX path semantics. Symlink races, time-of-check/time-of-use attacks and Windows paths are out of scope.

\paragraph{Threat model and policy.}
The attacker controls Flask request data: arguments, form fields, headers, cookies, the raw query string and the request path. Module-level configuration constants are trusted. One shared, versioned policy lists sources, sinks (file operations, restricted to their path argument) and a sanitizer (\texttt{secure\_filename}). It also lists the documented behaviour of three benchmark helper APIs, for example a request wrapper whose \texttt{get\_safe\_value} returns a constant. Every arm receives the same policy; no arm gets hidden knowledge.

\paragraph{Baseline: Semgrep taint mode.}
Semgrep~\cite{semgrep} matches code against patterns written in the target language. In taint mode a rule declares \emph{sources}, \emph{sinks}, \emph{sanitizers} and \emph{propagators}, and the engine reports any data-flow path from a source to a sink that does not pass through a sanitizer. Taint flows through assignments and, by default, through function calls: a call with a tainted argument yields a tainted result. Propagators add side-effect flows, such as \texttt{lst.append(x)} tainting \texttt{lst}. The Community Edition analyzes each function separately; following flows across functions and files is a feature of the commercial engine. We confirmed this on our own fixtures. Semgrep does not reason about values, so a variable overwritten by a constant, or assigned only on a branch whose condition is a false constant, still counts as tainted. It also does not interpret guard predicates.

The only community Python CWE-22 rule for Flask in the registry is pattern-based and matches only \texttt{open()}. Our subjects mostly use \texttt{pathlib} operators, \texttt{codecs.open} and existence checks. We therefore wrote our own taint rule from the policy, developed only on development data and then frozen. In pinned Semgrep 1.182.1 the Python engine does not propagate taint through string slices (\texttt{s[a:b]}). During development we tried propagators, rule options and taint labels without success, and kept this as a documented engine limitation. A second loss, values stored into a dictionary by subscript assignment and read back by key, appeared only in test-split failure analysis after the protocol was frozen. It was therefore not addressed.
